\chapter{Background tecnico}\label{chapter:background}
In questo capitolo vengono affrontate le conoscenze tecniche che supportano la comprensione della vulnerabilità \emph{format string} e vengono introdotti l'ambiente e gli strumenti utilizzati nel lavoro.


\section{Funzioni variadiche in C e specificatori di formato di printf}\label{sec:funcvar_format}
\begin{defn}[Funzione variadica]\label{def:funzione-variadica}
Nel linguaggio C, una \textbf{funzione variadica} è una funzione progettata per accettare un numero variabile di argomenti, potenzialmente di tipi diversi \cite{kr1988}.
\end{defn}

Una delle funzioni variadiche più note è \texttt{printf}, dichiarata nell'\emph{header} \texttt{<stdio.h>} della libreria standard del linguaggio C \cite{man3printf}, come mostrato nel Codice~\ref{lst:printf_decl}.
\begin{lstlisting}[
    language=C,
    numbers=none,
    gobble=4,
    caption={Dichiarazione di \texttt{printf}},
    label={lst:printf_decl}
]
    int printf(const char *fmt, ...);
\end{lstlisting}
L'ellissi \texttt{...} indica che il numero e il tipo di questi argomenti possono variare e che quindi si tratta di una funzione variadica. Questa ellissi può comparire solo alla fine di un elenco di argomenti.
Per le funzioni variadiche, il numero e il tipo degli argomenti che seguono l'ultimo parametro dichiarato non sono determinabili dalla sola dichiarazione della funzione. È quindi compito del programmatore definire una convenzione che il chiamante rispetti, per comunicare in qualche altro modo quanti argomenti sono stati passati. È possibile:
\begin{itemize}
    \item utilizzare una \emph{format string} che indichi il numero e il tipo dei parametri (ad es. la \emph{format string} di \texttt{printf});
    \item passare il numero come argomento fisso;
    \item terminare la lista degli argomenti aggiuntivi con un valore speciale (sentinella).
\end{itemize}
L'altro aspetto rilevante delle funzioni variadiche è che gli argomenti
catturati dall'ellissi non corrispondono ad alcun parametro nominato:
non è quindi possibile accedervi per nome, come si farebbe con un
parametro ordinario, ma serve un meccanismo diverso \cite{kr1988}.

Per inquadrare il funzionamento della \emph{format string} e il modo in cui le funzioni variadiche scorrono la lista degli argomenti aggiuntivi, guardiamo \texttt{minprintf} (Codice~\ref{lst:minprintf}), una versione molto semplificata di \texttt{printf} \cite{kr1988}.
\begin{lstlisting}[
    language=C,
    gobble=4,
    caption={Implementazione semplificata di \texttt{printf} (\texttt{minprintf})},
    label={lst:minprintf}
]
    void minprintf(char *fmt, ...)
    {
        va_list ap;
        va_start(ap, fmt);

        for (char *p = fmt; *p; p++) {
            if (*p != '%') {
                putchar(*p);
                continue;
            }
            switch (*++p) {
            case 'x':
                printf("%x", va_arg(ap, int));
                break;
            case 's':
                printf("%s", va_arg(ap, char *));
                break;
            // ...
            }
        }

        va_end(ap);
    }
\end{lstlisting}
Per gestire gli argomenti aggiuntivi, il linguaggio C ci fornisce l'header standard \texttt{<stdarg.h>}, contenente un insieme di definizioni di macro che stabiliscono come avanzare all'interno di una lista di argomenti. 
Il tipo \texttt{va\_list} viene utilizzato per dichiarare una variabile che farà riferimento, uno dopo l'altro, a ciascuno degli argomenti aggiuntivi anonimi; in \texttt{minprintf}, questa variabile è chiamata \texttt{ap}, che sta per ``argument pointer''. La macro \texttt{va\_start} inizializza \texttt{ap} facendolo puntare al primo argomento senza nome. Questa macro deve essere invocata una volta prima di poter utilizzare \texttt{ap}. È necessaria la presenza di almeno un argomento con nome; l'ultimo di questi argomenti viene utilizzato da \texttt{va\_start} come punto di partenza.

Ogni chiamata a \texttt{va\_arg} restituisce un argomento e fa avanzare \texttt{ap} a quello successivo; \texttt{va\_arg} utilizza il nome del tipo per determinare quale tipo di dato restituire e l'entità del passo da compiere per avanzare. Infine, \texttt{va\_end} si occupa di eseguire le eventuali operazioni di pulizia necessarie. Deve essere chiamata prima che la funzione termini \cite{kr1988}.

\begin{defn}[Specificatore di formato]\label{def:specificatore}
    Uno \textbf{specificatore di formato} è una sottosequenza della \emph{format string}, introdotta dal carattere \%, che stabilisce se e come un argomento aggiuntivo debba essere interpretato e restituito nell'output prodotto \cite{cplusplusprintf}.
\end{defn}

La funzione \texttt{minprintf} stampa sullo standard output (\texttt{stdout}) la stringa indicata da \texttt{fmt}. Qualora quest'ultima contenga degli specificatori di formato, gli argomenti successivi vengono opportunamente formattati e sostituiti al loro posto prima della stampa \cite{cplusplusprintf}. Tale meccanismo ricalca il tipico comportamento delle \emph{format string}.

Il principio di \texttt{printf} della libreria standard è identico; l'aspetto che cambia principalmente è la ricchezza dello specificatore di formato, che infatti ha la seguente struttura (Codice~\ref{lst:format_specifier_syntax}) \cite{man3printf,cplusplusprintf}:
\begin{lstlisting}[
    language={},
    numbers=none,
    gobble=4,
    caption={Struttura generale di uno specificatore di formato della famiglia \texttt{printf} (con l'estensione posizionale POSIX)},
    label={lst:format_specifier_syntax}
]
    %[n$][flag][larghezza][.precisione][lunghezza]specificatore
\end{lstlisting}
Per il nostro scopo sono rilevanti:
\begin{itemize}
    \item \emph{posizione} (\texttt{n\$}), permette di indicare esplicitamente quale argomento della lista debba essere formattato, anziché consumare implicitamente il successivo non ancora utilizzato (ad es. \texttt{\%3\$x} formatta come esadecimale il terzo argomento dopo \texttt{fmt}). Non fa parte dello standard ISO C, ma è un'estensione POSIX \cite{man3printf}; la sua disponibilità dipende quindi dalla libreria C effettivamente linkata, non dal linguaggio in sé.
    \item \emph{larghezza}, il numero minimo di caratteri da stampare. Se il valore da stampare è più corto di questo numero, il risultato viene riempito con spazi vuoti. Il valore non viene troncato, nemmeno se risulta più lungo.
    \item \emph{lunghezza}, modifica la dimensione del tipo di dato. Noi useremo:
        \begin{itemize}
            \item \texttt{hh} che per definizione è esattamente 1 byte di dimensione
            \item \texttt{h} che è praticamente sempre 2 byte di dimensione
        \end{itemize}

    \item \emph{specificatore}, definisce il tipo e l'interpretazione del relativo argomento. Ci interessano:
        \begin{itemize}
            \item \texttt{x}, stampa l'argomento intero esadecimale senza segno
            \item \texttt{p}, stampa l'argomento come indirizzo di memoria (puntatore)
            \item \texttt{n}, non viene stampato nulla, però il numero di caratteri scritti fino a quel momento viene memorizzato nella locazione di memoria puntata dall'argomento corrispondente
        \end{itemize}
\end{itemize}


\section{La calling convention x86\_64 e il confronto con x86}\label{sec:callconv}
La vulnerabilità \emph{format string}, la vediamo in particolare per l'architettura x86\_64, nota anche come AMD64. Ci interessano le convenzioni utilizzate da tale architettura durante la chiamata di funzione, in particolare dove vengono memorizzati gli argomenti passati alla funzione (ci servirà quando andremo a sfruttare la vulnerabilità).
Facciamo poi anche un rapido confronto con l'architettura x86 per mostrare l'impatto che il cambiamento di architettura può avere sulla tecnica di sfruttamento.

\subsection{System V AMD64 ABI: registri per il passaggio dei parametri}\label{subsec:x86_64}
\begin{defn}[ABI]\label{def:abi}
Un'\textbf{ABI} (\emph{Application Binary Interface}) definisce in che modo i diversi componenti di un codice binario possano interfacciarsi su un determinato sistema operativo e una determinata architettura \cite{silberschatz2019}.
\end{defn}

Tra i molti dettagli di basso livello specificati da un'ABI, quello che riguarda più da vicino la vulnerabilità trattata in questa tesi è la \emph{calling convention}, ovvero l'insieme di regole con cui il chiamante dispone gli argomenti di una funzione affinché il chiamato possa recuperarli correttamente. Per la piattaforma di riferimento (Linux, x86\_64), queste regole sono specificate dalla System V AMD64 ABI, che decide la collocazione di ciascun argomento --- in un registro o sullo stack --- sulla base di una classificazione.

Gli argomenti delle funzioni vengono suddivisi in diverse categorie, mappate sulle specifiche famiglie di registri dell'architettura AMD64. Quelli rilevanti ai fini della trattazione della nostra vulnerabilità appartengono alla classe \emph{INTEGER}, costituita dai tipi interi che possono essere contenuti all'interno di uno dei registri di uso generale (\emph{general purpose registers}).

I tipi base (con segno e senza segno) appartenenti alla classe \emph{INTEGER} sono: \texttt{\_Bool}, \texttt{char}, \texttt{short}, \texttt{int}, \texttt{long}, \texttt{long long} e puntatori.
La dimensione di ciascun argomento viene arrotondata per eccesso a otto byte. Di conseguenza, lo stack risulta sempre allineato a otto byte \cite{psabi2025}.

Per gli argomenti di classe \emph{INTEGER}, il passaggio avviene assegnandoli da sinistra verso destra al primo registro disponibile nella sequenza: \texttt{\%rdi}, \texttt{\%rsi}, \texttt{\%rdx}, \texttt{\%rcx}, \texttt{\%r8}, \texttt{\%r9} (Tabella~\ref{tab:registri_integer}).
Completata questa assegnazione ai registri, gli argomenti rimanenti (destinati alla memoria) vengono inseriti nello stack in ordine inverso, ovvero da destra verso sinistra \cite{psabi2025}.

\begin{table}[ht]
    \centering
    \caption{Registro assegnato a ciascuna posizione per gli argomenti di classe \emph{INTEGER}}
    \label{tab:registri_integer}
    \begin{tabular}{cc}
        \toprule
        Posizione argomento & Registro \\
        \midrule
        1 & \texttt{\%rdi} \\
        2 & \texttt{\%rsi} \\
        3 & \texttt{\%rdx} \\
        4 & \texttt{\%rcx} \\
        5 & \texttt{\%r8}  \\
        6 & \texttt{\%r9}  \\
        \midrule
        $\geq 7$ & stack \\
        \bottomrule
    \end{tabular}
\end{table}

\subsection{Il caso delle funzioni variadiche}\label{subsec:funvar_cc}
Le funzioni standard, a parametri fissi, applicano la \emph{calling convention} x86\_64 senza difficoltà, poiché conoscono in anticipo numero e tipo degli argomenti che vengono loro passati, e sanno dunque già dove reperirli.

Le funzioni variadiche, invece, non sanno in anticipo quanti argomenti riceveranno, né se si troveranno nei registri o sullo stack.
Come visto in \ref{sec:funcvar_format}, le funzioni variadiche gestiscono gli argomenti aggiuntivi tramite l'\emph{header} \texttt{<stdarg.h>}, come se fossero una lista. Però non abbiamo specificato l'implementazione di questa lista, che varia da architettura a architettura poiché ci sono convenzioni differenti da rispettare. Vediamo dunque come \texttt{<stdarg.h>} sia implementato concretamente sull'architettura x86\_64.

In x86\_64, i primi argomenti interi vengono passati nei sei registri nella Tabella~\ref{tab:registri_integer}. Il problema è che \texttt{<stdarg.h>} usa dei puntatori strutturati per accedere agli argomenti e non possiamo accedere tramite puntatori a dei registri come invece facciamo per gli argomenti nello stack.
La soluzione è stata quella di mettere nel prologo delle funzioni variadiche del codice che vada a salvare il contenuto dei registri usati per il passaggio degli argomenti in un'area di memoria \ref{def:reg_save_area}.

\begin{defn}[Register Save Area]\label{def:reg_save_area}
    La \textbf{Register Save Area} è un'area di memoria a dimensione fissa (176 byte), allocata nello \emph{stack frame} di una funzione variadica che invoca \texttt{va\_start}.
    Ogni registro usato per il passaggio dei parametri vi occupa un offset fisso (Figura~\ref{tab:register_save_area}), così che \texttt{va\_arg} possa accedere agli argomenti arrivati nei registri con la stessa aritmetica dei puntatori usata per quelli sullo stack.
\end{defn}

\begin{table}[ht]
    \centering
    \caption{Offset dei registri argomento nella \emph{register save area}}
    \label{tab:register_save_area}
    \begin{tabular}{cc}
        \toprule
        Registro & Offset (byte) \\
        \midrule
        \texttt{\%rdi} & 0 \\
        \texttt{\%rsi} & 8 \\
        \texttt{\%rdx} & 16 \\
        \texttt{\%rcx} & 24 \\
        \texttt{\%r8}  & 32 \\
        \texttt{\%r9}  & 40 \\
        \midrule
        \texttt{\%xmm0--\%xmm7} & 48--175 \\
        \bottomrule
    \end{tabular}
\end{table}

Devono essere salvati solo i registri che potrebbero essere utilizzati per passare gli argomenti. 
Una funzione variadica chiamata non sa quanti registri verranno popolati \todo{nota a piè pagina per i registri vettoriali e \%al che fa limite superiore}, di conseguenza, per non perdere potenziali dati validi, il prologo è obbligato a salvare sempre tutti e 6 i registri interi destinati al passaggio dei parametri all'interno della register save area sullo stack.

Il tipo \texttt{va\_list} (Codice~\ref{lst:va_list}) è un array di un solo elemento di tipo \texttt{struct}, contenente le informazioni necessarie per permettere alla macro \texttt{va\_arg} di estrarre l'argomento successivo correttamente. È compito della macro \texttt{va\_start} di inizializzare questa \texttt{struct}.

\begin{lstlisting}[
    language=C,
    numbers=none,
    gobble=4,
    caption={\texttt{va\_list} type declaration},
    label={lst:va_list}
]
    typedef struct {
        unsigned int gp_offset;
        unsigned int fp_offset;
        void *overflow_arg_area;
        void *reg_save_area;
    } va_list[1];
\end{lstlisting}

Vediamo le componenti di \texttt{va\_list}:
\begin{itemize}
    \item \texttt{reg\_save\_area}: puntatore all'inizio della \emph{register save area}.

    \item \texttt{gp\_offset}: contiene l'offset in byte, a partire da \texttt{reg\_save\_area}, del punto in cui è salvato il successivo registro \emph{general-purpose} disponibile. Nel caso in cui tutti i registri di argomento siano esauriti, viene impostato al valore 48 ($6 \times 8$).

    \item \texttt{fp\_offset}: contiene l'offset in byte, a partire da \texttt{reg\_save\_area}, del punto in cui è salvato il successivo registro \emph{floating-point} disponibile. Nel caso in cui tutti i registri di argomento siano esauriti, viene impostato al valore 176 ($6 \times 8 + 8 \times 16$).

    \item \texttt{overflow\_arg\_area}: puntatore nello stack "vero", fuori dalla register save area: punta al primo argomento passato lì (se esiste) e avanza a ogni lettura successiva. È dove \texttt{va\_arg} finisce per leggere una volta che \texttt{gp\_offset} o \texttt{fp\_offset} hanno raggiunto il loro valore massimo.
\end{itemize}

Come detto più volte, la macro \texttt{va\_arg} ci permette di estrarre gli argomenti aggiuntivi di una funzione variadica correttamente.
A ogni chiamata di \texttt{va\_arg(ap, type)} (ap è di tipo \texttt{va\_list}) controlla se \texttt{type} è idoneo a essere passato tramite registri. In caso
\begin{itemize}
    \item \emph{positivo}, viene calcolato quanti registri sono necessari per contenere \texttt{type} e se i registri
    \begin{itemize}
        \item \emph{sono disponibili}, l'argomento viene estratto dai relativi registri in base agli offset correnti.
        \item \emph{non sono disponibili}, l'argomento viene estratto dallo stack, in base agli offset correnti
    \end{itemize}

    \item \emph{negativo}, l'argomento viene estratto dallo stack, in base agli offset correnti.
\end{itemize}
Alla fine di ogni chiamata, vengono aggiornati i relativi offset in \texttt{va\_list} in base ai registri utilizzati o allo spazio consumato nello stack.

Vediamo ora come questi meccanismi si applichino al caso specifico \texttt{printf}.
La funzione variadica \texttt{printf} possiede la \emph{format string} come unico parametro fisso. Tale argomento occuperà quindi il primo registro \emph{general-purpose} disponibile, ovvero \texttt{\%rdi}. Gli eventuali argomenti aggiuntivi, occuperanno i successivi registri o lo stack, in base al loro tipo e numero.
Se al chiamante passata esclusivamente la \emph{format string} senza ulteriori argomenti, il prologo della funzione provvederà comunque a salvare il contenuto (potenzialmente indefinito) dei restanti cinque registri \emph{general-purpose} all'interno della \emph{register save area}.
Qualora la \emph{format string} contenesse degli specificatori (ad es. \texttt{\%x} o \texttt{\%p}), la libreria richiamerebbe iterativamente \texttt{va\_arg}, andando a leggere il contenuto di quei registri che, di fatto, non corrispondono ad alcun argomento reale. Per questo motivo, ai fini di questa trattazione, possiamo riferirci a questi valori con il termine informale di ``argomenti fantasma'' (ovvero dati residui di precedenti operazioni trattati erroneamente come parametri validi).
Se, inoltre, la stringa contenesse un numero di specificatori tale da esaurire i registri disponibili, \texttt{va\_arg} proseguirebbe la lettura, andando a prelevare ulteriori ``argomenti fantasma'',direttamente dallo stack.
Vedremo nel Capitolo 2 come questo comportamento permetta a un attaccante di ottenere la lettura e la scrittura di locazioni di memoria arbitrarie, gettando le basi teoriche per lo sfruttamento delle vulnerabilità di tipo \emph{Format String}.
